\CWHeader{Лабораторная работа \textnumero 2\\ \enquote{Поисковый робот}}

Необходимо написать поисковый робот --- компоненту обкачки документов.

\begin{itemize}
    \item Единственным аргументом роботу подаётся путь до yaml-конфига,
          содержащего данные для базы данных в секции \texttt{db}, данные для
          робота в секции \texttt{logic} (задержка между обкачкой страницы) и
          любые другие данные, необходимые для реализации логики робота.
    \item Сохранять в базе данных (например, MongoDB) документы с полями:
          нормализованный url, \enquote{сырой} html-текст документа, название
          источника, дата обкачки в формате Unix timestamp.
    \item Робот можно остановить в любой момент, и при повторном запуске он
          должен начать с того документа, на котором остановился.
    \item Периодически робот должен уметь переобкачивать документы, которые уже
          есть в базе, но только если они изменились.
\end{itemize}

\pagebreak
